\begin{trapbox}

	\begin{description}[style=nextline, leftmargin=2.8em, labelsep=0.5em, itemsep=0.35em]
		\item[\textbf{\textcolor{CTrap}{易错点一(焦点方向)}}]
			误认为双曲线焦点三角形面积公式只适用于焦点在x轴的情形，当方程为$\frac{y^2}{a^2}-\frac{x^2}{b^2}=1$时不敢直接使用$S=b^2\cot\frac{\theta}{2}$。
		\item[\textbf{\textcolor{CTrap}{正确理解(形式不变)：}}]
			对于两种标准形式，公式中的b始终是虚半轴（即减号后的分母），公式完全一样。
		\item[\textbf{\textcolor{CTrap}{错因分析(忽视推导)：}}]
			推导过程中只用了$b^2=c^2-a^2$，并未假定焦点位置，故公式与焦点取向无关。
	\end{description}

	\medskip
	\begin{description}[style=nextline, leftmargin=2.8em, labelsep=0.5em, itemsep=0.35em]
		\item[\textbf{\textcolor{CTrap}{易错点二(忽视退化)}}]
			应用公式时没有验证点P不是实轴端点，直接代入公式。
		\item[\textbf{\textcolor{CTrap}{正确理解(条件限制)：}}]
			当P为顶点时，三角形退化为线段，不能用此公式；公式推导中假设了$r_1\neq r_2$且$\theta\neq0,\pi$。
		\item[\textbf{\textcolor{CTrap}{错因分析(无视条件)：}}]
			解题时只关注面积形式，忽略前提条件。
	\end{description}

	\medskip
	\begin{description}[style=nextline, leftmargin=2.8em, labelsep=0.5em, itemsep=0.35em]
		\item[\textbf{\textcolor{CTrap}{易错点三(参数误判)}}]
			在双曲线方程中，误将实半轴平方$a^2$当作公式中的$b^2$代入计算。
		\item[\textbf{\textcolor{CTrap}{正确理解(明确参数)：}}]
			公式中的b是虚半轴，对应方程中减号后面分母的平方根。
		\item[\textbf{\textcolor{CTrap}{错因分析(概念不清)：}}]
			未区分实轴与虚轴在公式中的角色。
	\end{description}

\end{trapbox}
